\begin{abstract}
Robots that revisit changing environments need to keep useful geometry from
earlier runs without preserving object positions that are no longer valid.
Starting every run from scratch loses coverage, while simply accumulating maps
leaves stale geometry behind. We present a method for updating a dense dynamic
map across separate mapping runs. Built on Khronos, the method stores successive
object states and background geometry, then updates them with new RGB-D
observations. We project new depth measurements onto previous surfaces to
distinguish agreement, occlusion by a closer surface, and evidence that a
previous surface is absent. Evidence is accumulated before removing an old
object state, and reconstructed surfaces are compared to decide whether a new
observation refines the previous state or represents a new one. On three
real-room runs, F1 for detecting object changes between runs is 96.8\% and
100.0\%, compared with 64.0\% and 37.5\% for Panoptic Multi-TSDFs, while
obsolete surface left where objects used to be drops from 57--60\% to 7--11\%.
On controlled synthetic data, the method reaches 97.2--97.7\% geometry F1 and
detects 29 of 32 changes between runs without a false event. The method
preserves previously observed geometry while removing obsolete object states
across repeated mapping runs.
\end{abstract}
